% Shared pleading-paper preamble for CGC-25-631802 TAC-RESPONSE court PDFs.
\documentclass{article}
\usepackage[utf8]{inputenc}
\usepackage{graphicx}
\usepackage{geometry}
    \geometry{top=1in, bottom=2.2in, left=1.0in, right=1.0in, textheight=336pt}
    \setlength{\footskip}{55pt}
\usepackage{ulem}
\usepackage{fancyhdr}
    \renewcommand{\headrulewidth}{0pt}
    \renewcommand{\footrulewidth}{0pt}
\usepackage{parskip}
    \setlength{\parskip}{0mm}
\usepackage{quoting}
    \quotingsetup{leftmargin=0.5in, rightmargin=0.5in, vskip=-1.5mm}
    \makeatletter
        \g@addto@macro\quoting\singlespacing
        \g@addto@macro\quoting{\vspace{-2mm}}
    \makeatother
\usepackage{mathptmx}
\usepackage{amssymb}
\renewcommand{\rmdefault}{ptm}
\renewcommand{\normalsize}{\fontsize{12}{12}\selectfont}
\newcommand*{\ptm}{\fontfamily{ptm} \selectfont \fontsize{12}{12}\selectfont}
\usepackage{setspace}
    \doublespacing
\newcommand{\tab}{\hspace*{.0in}}
\newenvironment{tightcenter}{%
    \setlength\topsep{0pt}
    \setlength\parskip{0pt}
    \begin{center}
}{%
    \end{center}
}
\usepackage{tikz}
\usetikzlibrary{calc}
\AddToHook{shipout/background}{%
    \begin{tikzpicture}[remember picture, overlay]
        \draw[line width=0.5pt] ($(current page.north west) + (0.75in,0)$) -- ($(current page.south west) + (0.75in,0)$);
        \draw[line width=0.5pt] ($(current page.north west) + (0.80in,0)$) -- ($(current page.south west) + (0.80in,0)$);
    \end{tikzpicture}%
}
\usepackage[left, pagewise]{lineno}
\setlength{\linenumbersep}{0.35in}
\renewcommand{\linenumberfont}{\normalfont\small}
\usepackage{titlesec}
\usepackage[hidelinks]{hyperref}
\usepackage{bookmark}
\usepackage{textcomp}
\usepackage{longtable}
\usepackage{booktabs}
\usepackage{array}
\usepackage{multirow}
\usepackage{calc}
% Pandoc pipe-table column widths use \real{#}; map to pgf math (pgf loaded via tikz).
\newcommand{\real}[1]{\pgfmathparse{#1}\pgfmathresult}
\titleformat{\section}{\fontsize{15}{15}\selectfont\normalfont\bfseries\uppercase}{}{0em}{}
\titleformat{\subsection}{\fontsize{15}{15}\selectfont\normalfont\bfseries\uppercase}{}{0em}{}
\titleformat{\subsubsection}{\normalfont\bfseries}{}{0em}{}
\titlespacing*{\section}{0pt}{0pt}{0pt}
\titlespacing*{\subsection}{0pt}{0pt}{0pt}
\titlespacing*{\subsubsection}{0pt}{0pt}{0pt}
\newcommand{\calrules}[1]{Cal. Rules of Court, rule #1}
\newcommand{\ccp}[1]{Code Civ. Proc., \textsection~#1}
\newcommand{\evidcode}[1]{Evid. Code, \textsection~#1}
\newcommand{\vehcode}[1]{Veh. Code, \textsection~#1}
\providecommand{\tightlist}{%
  \setlength{\itemsep}{0pt}\setlength{\parskip}{0pt}}
% --- TOC / TOA hyperlink infrastructure ---
\newcommand{\tacctoclink}[2]{\hyperref[toc:#1]{#2}}
% In-text citation link to authority label (target is usually TABLE OF AUTHORITIES row).
\newcommand{\citecase}[2]{\textit{\hyperref[auth:#1]{#2}}}
\newcommand{\citestatute}[2]{\hyperref[auth:#1]{#2}}
\newcommand{\toclink}[2]{\hyperref[toc:#1]{#2}}
% Authority entries in TABLE OF AUTHORITIES (#1 label slug, #2 italic case name or title, #3 citation tail)
\newcommand{\caseentry}[3]{\par\noindent\hangindent=2em\hangafter=1%
  \phantomsection\label{auth:#1}\textit{#2}, #3\par}
\newcommand{\statuteentry}[3]{\par\noindent\hangindent=2em\hangafter=1%
  \phantomsection\label{auth:#1}#2\ #3\par}
\newcommand{\ruleentry}[2]{\par\noindent\hangindent=2em\hangafter=1%
  \phantomsection\label{auth:#1}#2\par}
\makeatletter
\newlength{\CourtTOCmarkwd}
\newlength{\CourtTOCpagewd}
\newlength{\CourtTOCtitlewd}
\setlength{\CourtTOCmarkwd}{2.8em}
\setlength{\CourtTOCpagewd}{4.2em}
\newcommand{\CourtTOCRow}[3]{%
  \begingroup
  \setlength{\CourtTOCtitlewd}{\dimexpr\linewidth - \CourtTOCmarkwd - \CourtTOCpagewd - 0.6em\relax}%
  \par\nobreak\noindent
  \begin{minipage}[b]{\CourtTOCmarkwd}\raggedright\strut #1\strut\end{minipage}\hspace{0.3em}%
  \begin{minipage}[b]{\CourtTOCtitlewd}\raggedright\strut #2\strut\end{minipage}\hspace{0.3em}%
  \begin{minipage}[b]{\CourtTOCpagewd}\raggedleft\strut #3\strut\end{minipage}\par
  \endgroup
  \vspace{1pt}%
}
\newcommand{\CourtTOCSubitem}[2]{\CourtTOCRow{}{#1}{#2}}
\newcommand{\CourtTOCSubsubitem}[2]{\CourtTOCRow{}{\quad\quad #1}{#2}}
\makeatother

\begin{document}
\nolinenumbers
\begin{singlespace}
Franciscus Dylan Rosario, \\
7624 Melody \\
Rohnert Park, CA 94928 \\
E: soltrinox@gmail.com \\
Plaintiff, In Pro Per \\
\end{singlespace}
\vspace*{9mm}
\begin{tightcenter}
\textbf{SUPERIOR COURT OF THE STATE OF CALIFORNIA} \\
\textbf{COUNTY OF SAN FRANCISCO}
\end{tightcenter}
\vspace*{2.25mm}
\nolinenumbers
\begin{minipage}[t]{3in}
    \raggedright
    \textbf{Franciscus Dylan Rosario,} \\ \textbf{PLAINTIFF}, \\
    \hspace{1cm} v. \\
    \small
    CSAA Insurance Exchange; \\
    and DOES 1 through 50, inclusive; \\
    \normalsize
    \textbf{DEFENDANTS}
    \makebox[3in]{\hrulefill}
\end{minipage}%
\hspace{3mm}
\begin{minipage}[t]{0.5pt}
    \vspace{0pt}
    \rule{0.5pt}{2.0in}
\end{minipage}
\hspace{3mm}
\begin{minipage}[t]{3in}
    \raggedright
    \noindent \textbf{Case No.: CGC-25-631802} \\
    ~\\
    \textbf{THIRD AMENDED COMPLAINT} \\
    \textbf{FOR FRAUD AND UNFAIR BUSINESS PRACTICES}
\end{minipage}
\vspace*{6mm}
\linenumbers
\begin{center}\textbf{TABLE OF CONTENTS}\end{center}
\vspace{2mm}
\CourtTOCRow{}{\tacctoclink{section-i-directduty-gravamen}{SECTION I: DIRECT-DUTY GRAVAMEN}}{\pageref{toc:section-i-directduty-gravamen}}
\CourtTOCRow{}{\tacctoclink{section-ii-prelitigation-privilege-carveout}{SECTION II: PRE-LITIGATION PRIVILEGE CARVEOUT}}{\pageref{toc:section-ii-prelitigation-privilege-carveout}}
\CourtTOCRow{}{\tacctoclink{section-iii-attorneys-fees-cure}{SECTION III: ATTORNEY'S FEES CURE}}{\pageref{toc:section-iii-attorneys-fees-cure}}
\CourtTOCRow{}{\tacctoclink{section-iv-punitive-damages-specificity}{SECTION IV: PUNITIVE DAMAGES SPECIFICITY}}{\pageref{toc:section-iv-punitive-damages-specificity}}
\CourtTOCRow{}{\tacctoclink{section-v-damages}{SECTION V: DAMAGES}}{\pageref{toc:section-v-damages}}
\CourtTOCRow{}{\tacctoclink{section-vi-primary-rights-reservation}{SECTION VI: PRIMARY RIGHTS / RESERVATION}}{\pageref{toc:section-vi-primary-rights-reservation}}
\CourtTOCRow{}{\tacctoclink{section-vii-title-10-investigation-standard}{SECTION VII: TITLE 10 INVESTIGATION STANDARD}}{\pageref{toc:section-vii-title-10-investigation-standard}}
\CourtTOCRow{}{\tacctoclink{section-viii-delayed-discovery}{SECTION VIII: DELAYED DISCOVERY}}{\pageref{toc:section-viii-delayed-discovery}}
\CourtTOCRow{}{\tacctoclink{causes-of-action}{CAUSES OF ACTION}}{\pageref{toc:causes-of-action}}
\CourtTOCRow{}{\tacctoclink{prayer-for-relief}{PRAYER FOR RELIEF}}{\pageref{toc:prayer-for-relief}}
\CourtTOCRow{}{\tacctoclink{demand-for-jury-trial}{DEMAND FOR JURY TRIAL}}{\pageref{toc:demand-for-jury-trial}}
\CourtTOCRow{}{\tacctoclink{verification}{VERIFICATION}}{\pageref{toc:verification}}
\vspace{4mm}
\phantomsection\label{toc:section-i-directduty-gravamen}
\section{SECTION I: DIRECT-DUTY GRAVAMEN}\label{section-i-direct-duty-gravamen}

\subsection{Title 10 CCR Direct Claimant Obligations}\label{title-10-ccr-direct-claimant-obligations}

\begin{enumerate}
\def\labelenumi{\arabic{enumi}.}
\item
  This action arises from Defendant CSAA Insurance Exchange's fraudulent handling of Plaintiff's claim as a direct third-party claimant. The actionable wrong is CSAA's non-communicative institutional failure to investigate the claim prior to denying it, and the active concealment of that failure through a denial letter representing the investigation as ``concluded.''
\item
  On \textbf{February 4, 2021}, Plaintiff was involved in a motor-vehicle collision in San Francisco County with CSAA's insured, Subhi Abdelhalim.
\item
  On \textbf{February 25, 2021}, CSAA sent Plaintiff a written denial letter concerning \textbf{Claim No.~1004-09-1054}, signed by \textbf{Sarah Rash, Claims Representative}. The denial letter stated, in pertinent part:
\end{enumerate}

\begin{quote}
``We have concluded our investigation of your claim. Based on our investigation, and after carefully reviewing the facts and circumstances of your claim, we have concluded that our insured is not liable\ldots{} Our decision is based on the facts currently available.''
\end{quote}

\begin{enumerate}
\def\labelenumi{\arabic{enumi}.}
\setcounter{enumi}{3}
\tightlist
\item
  This representation was false. At the time of the denial, CSAA had not conducted the thorough investigation required by Title 10 CCR \textsection{}\textsection{} 2695.4 and 2695.7.
\end{enumerate}

\subsection{\textsection{}I.5 --- Direct-Addressee Theory (Primary Shaolian Rebuttal)}\label{i.5-direct-addressee-theory-primary-shaolian-rebuttal}

\begin{enumerate}
\def\labelenumi{\arabic{enumi}.}
\setcounter{enumi}{4}
\item
  Plaintiff asserts standing as the \textbf{direct addressee} of a specific misrepresentation communicated by CSAA to Plaintiff individually. This is not a claim for policy benefits, not a claim for breach of an insurance contract by an insured, and not a third-party bad-faith claim derived from the insured's contractual relationship with CSAA. It is a common law fraud claim by the person to whom the representation was made, on whom it was intended to operate, and who suffered individual and direct economic harm as a result.
\item
  By sending the February 25, 2021 denial letter directly to Plaintiff as an unrepresented adverse claimant, rather than communicating through an insured or counsel, CSAA voluntarily undertook a direct communicative relationship with Plaintiff. That voluntary direct communication, combined with CSAA's knowledge that Plaintiff was unrepresented and had no independent means to verify the scope of CSAA's investigation, establishes the direct duty context for common law fraud liability.
\end{enumerate}

6A. \textbf{Agency and Ratification Under Civil Code \textsection{} 2338.} Civil Code \textsection{} 2338 provides: ``Unless required by or under the authority of law to employ that particular agent, a principal is responsible to third persons for the negligence of his agent in the transaction of the business of the agency, including wrongful acts committed by such agent in and as a part of the transaction of such business, and for his willful omission to fulfill the obligations of the principal.'' Under \textsection{} 2338, CSAA Insurance Exchange is directly liable for the wrongful acts and omissions of Sarah Rash, the claims department personnel who reviewed and approved the February 25, 2021 denial, and any retained claims professionals or counsel acting within the scope of the claims-handling agency. Plaintiff further alleges that CSAA ratified the February 25, 2021 denial and its underlying non-investigation by (a) failing to correct the denial after the Title 10 CCR \textsection{}\textsection{} 2695.4 and 2695.7 violations became apparent; (b) maintaining the denial position through successive insurance-coverage communications and, later, through litigation defense of its insured; and (c) failing to repudiate the ``concluded investigation'' representation after the January 17, 2025 / February 18, 2025 / April 16, 2025 record events described in Section VIII. Under \emph{Doctors' Company v. Superior Court} (1989) 49 Cal.3d 39, an insurer that acts outside the ordinary insurer role by directing, approving, or ratifying wrongful claims-handling conduct may be held directly liable for that conduct. This paragraph forecloses the anticipated defense that ``CSAA itself did not act'' by pleading principal-agent and ratification liability in express Civil Code \textsection{} 2338 terms.

\subsection{\textsection{}I.6 --- Gravamen Disclaimer (Non-Communicative Institutional Failure)}\label{i.6-gravamen-disclaimer-non-communicative-institutional-failure}

\begin{enumerate}
\def\labelenumi{\arabic{enumi}.}
\setcounter{enumi}{6}
\item
  The gravamen of this action is CSAA's institutional ``deny-delay-defend'' failure to investigate, not the act of writing a denial letter. The February 25, 2021 denial letter is pleaded as \textbf{evidence} of the underlying non-communicative failure, not as the sole injury-producing act.
\item
  Among the foundational incident materials CSAA failed to obtain, review, or truthfully account for before the February 25, 2021 denial were the following contemporaneous San Francisco first-responder records, each of which was available from the San Francisco Police Department and/or the San Francisco Department of Emergency Management within days of the February 4, 2021 collision and was obtainable by routine written request under Gov.~Code \textsection{} 6250 et seq. (now Gov.~Code \textsection{}\textsection{} 7920.000 et seq.) and Penal Code \textsection{} 832.7(b):

  \begin{enumerate}
  \def\labelenumii{(\alph{enumii})}
  \item
    The \textbf{911 audio recording} and associated \textbf{Computer-Aided Dispatch (CAD) event record} generated on February 4, 2021 in connection with \textbf{SFPD Incident Report No.~210083458} (on information and belief; exact SFPD incident and 911-CAD event numbers to be conformed to proof), in which CSAA's insured made statements acknowledging that he struck Plaintiff with his vehicle;
  \item
    \textbf{Plaintiff's contemporaneous report to 911} emergency services, made on February 4, 2021 from the collision scene, that the driver had struck him; and
  \item
    \textbf{Body-worn camera (BWC) footage} from the San Francisco Police Department officers who responded to the February 4, 2021 collision scene (responding officers' names and star numbers in the possession of SFPD and available through a routine Penal Code \textsection{} 832.7(b) / Pitchess-adjacent public-records request), in which a third-party eyewitness (\textbf{Justin Rosemond}) independently corroborated that the driver struck Plaintiff.
  \end{enumerate}
\item
  Taken together, these three sources constituted virtually uncontested evidence of liability---confession, victim report, and eyewitness corroboration. CSAA denied a claim supported by all three.
\end{enumerate}

\subsection{UCL Standing Under Business and Professions Code \textsection{} 17200}\label{ucl-standing-under-business-and-professions-code-17200}

\begin{enumerate}
\def\labelenumi{\arabic{enumi}.}
\setcounter{enumi}{9}
\tightlist
\item
  Plaintiff has standing to bring a claim under Business and Professions Code \textsection{} 17200 et seq. (the Unfair Competition Law) because Plaintiff has suffered injury in fact and has lost money or property as a result of Defendant's unfair, fraudulent, and unlawful business practices.
\end{enumerate}

\subsection{Express Moradi-Shalal Carve-Out}\label{express-moradi-shalal-carve-out}

\begin{enumerate}
\def\labelenumi{\arabic{enumi}.}
\setcounter{enumi}{10}
\item
  This action does not assert a private cause of action under Insurance Code \textsection{} 790.03, which is precluded by \emph{Moradi-Shalal v. Fireman's Fund Insurance Companies} (1988) 46 Cal.3d 287.
\item
  Instead, Plaintiff asserts:
\end{enumerate}

\begin{enumerate}
\def\labelenumi{(\alph{enumi})}
\item
  \textbf{Common law fraud and deceit} under Civil Code \textsection{}\textsection{} 1709, 1710;
\item
  \textbf{Fraud --- Concealment / Half-Truth}; and
\item
  \textbf{Unfair Competition Law (UCL)} claims under Business and Professions Code \textsection{} 17200, which permits predicate violations of Insurance Code \textsection{} 790.03(h) as evidence of ``unlawful'' conduct.
\end{enumerate}

\begin{enumerate}
\def\labelenumi{\arabic{enumi}.}
\setcounter{enumi}{12}
\tightlist
\item
  Plaintiff's UCL claim is not barred by \emph{Moradi-Shalal} because the UCL provides an independent statutory cause of action that may be predicated on violations of Insurance Code \textsection{} 790.03. See \emph{State Farm Fire \& Casualty Co.~v. Superior Court} (1996) 45 Cal.App.4th 1093, 1103--04.
\end{enumerate}

\begin{center}\rule{0.5\linewidth}{0.5pt}\end{center}

\phantomsection\label{toc:section-ii-prelitigation-privilege-carveout}
\section{SECTION II: PRE-LITIGATION PRIVILEGE CARVEOUT}\label{section-ii-pre-litigation-privilege-carveout}

\subsection{Denial Letter as Pre-Litigation Communication}\label{denial-letter-as-pre-litigation-communication}

\begin{enumerate}
\def\labelenumi{\arabic{enumi}.}
\setcounter{enumi}{13}
\item
  The February 25, 2021 denial letter was a pre-litigation claim-handling communication, not a litigation communication protected by Civil Code \textsection{} 47(b).
\item
  Under \emph{Action Apartment Association, Inc.~v. City of Santa Monica} (2007) 41 Cal.4th 1232, 1248--49, the litigation privilege applies only to communications made in judicial or quasi-judicial proceedings. An insurer's denial of a claim, made in the ordinary course of claims handling, is not a litigation communication.
\item
  At the time of the February 25, 2021 denial:
\end{enumerate}

\begin{enumerate}
\def\labelenumi{(\alph{enumi})}
\item
  No litigation was pending or contemplated;
\item
  No attorney had been retained by either party;
\item
  The denial was issued in the ordinary course of CSAA's claims handling process; and
\item
  The denial was issued just 21 days after the collision---before any meaningful investigation could have been completed.
\end{enumerate}

16A. \textbf{Action Apartment ``Contemplation'' Factors --- Affirmative Pleading of Each Prong's Absence.} Under \emph{Action Apartment Ass'n, Inc.~v. City of Santa Monica} (2007) 41 Cal.4th 1232, 1248--49, a pre-litigation communication falls within Civil Code \textsection{} 47(b) only when it is made ``in good faith {[}and{]} in serious contemplation'' of litigation that is ``imminently anticipated.'' That test has three operative factors, each of which must be affirmatively present at the time of the communication: (i) litigation must be seriously proposed by the communicating party and in good faith contemplated as a meaningful possibility, not a remote or speculative one; (ii) the communication must have some functional connection to that contemplated litigation --- e.g., preservation of evidence, a litigation-hold directive, a formal demand, or a tolling communication; and (iii) the communicating party must itself be a contemplated party to that litigation or its authorized representative. Plaintiff affirmatively alleges, as of February 25, 2021, that \textbf{none} of the three Action Apartment factors was present in CSAA's favor: \textbf{(a)} no litigation was pending, filed, threatened in writing, or seriously contemplated by CSAA in good faith at the time of the denial letter; \textbf{(b)} no counsel had been retained by CSAA or by Plaintiff in connection with the claim, no attorney had appeared on behalf of either party, and Plaintiff was an unrepresented direct third-party claimant; \textbf{(c)} no pre-suit demand letter had been exchanged, no tolling agreement had been proposed or executed, no litigation-hold or document-preservation letter had been sent, and no litigation reserve had been set by CSAA for the claim as a litigation matter (as opposed to a routine claims reserve). Because each Action Apartment factor fails, Civil Code \textsection{} 47(b) does not attach to the February 25, 2021 denial.

\subsection{Ribas Distinction --- Fraud, Not IIED}\label{ribas-distinction-fraud-not-iied}

\begin{enumerate}
\def\labelenumi{\arabic{enumi}.}
\setcounter{enumi}{16}
\item
  To the extent Defendant cites \emph{Ribas v. Clark} (1985) 38 Cal.3d 355 for a bar on emotional distress claims against insurers, Plaintiff notes that \emph{Ribas} addresses intentional infliction of emotional distress (IIED) claims, not fraud claims.
\item
  Plaintiff does not assert a standalone IIED cause of action. Emotional distress damages are sought only as parasitic damages accompanying the fraud tort.
\item
  \emph{Ribas} does not bar fraud claims where, as here, the gravamen is the misrepresentation of investigative completeness, not outrageous conduct intended to cause emotional distress.
\end{enumerate}

\subsection{Non-Communicative Conduct --- Failure to Investigate}\label{non-communicative-conduct-failure-to-investigate}

\begin{enumerate}
\def\labelenumi{\arabic{enumi}.}
\setcounter{enumi}{19}
\tightlist
\item
  Defendant's failure to obtain and review foundational incident records constitutes non-communicative conduct that is not protected by Civil Code \textsection{} 47(b):
\end{enumerate}

\begin{enumerate}
\def\labelenumi{(\alph{enumi})}
\item
  \textbf{911 Audio:} Defendant failed to obtain the 911 call recording containing the driver's own acknowledgment;
\item
  \textbf{Body-Worn Camera Footage:} Defendant failed to obtain police BWC footage containing independent eyewitness corroboration;
\item
  \textbf{CAD Records:} Defendant failed to obtain Computer-Aided Dispatch records documenting the emergency response.
\end{enumerate}

\begin{enumerate}
\def\labelenumi{\arabic{enumi}.}
\setcounter{enumi}{20}
\tightlist
\item
  This failure to investigate is non-communicative because it consists of omissions and failures to act, not statements or publications. \emph{Rusheen v. Cohen} (2006) 37 Cal.4th 1048 (privilege does not apply to non-communicative conduct).
\end{enumerate}

21A. \textbf{Flatley Illegality Carve-Out.} Independently of the Action Apartment pre-litigation analysis above, Civil Code \textsection{} 47(b) does not protect conduct that is illegal as a matter of law. \emph{Flatley v. Mauro} (2006) 39 Cal.4th 299, 317--21, holds that ``where a defendant brings a motion to strike under section 425.16 based on a claim that the plaintiff's action arises from activity by the defendant in furtherance of the defendant's exercise of protected speech or petition rights, but either the defendant concedes, or the evidence conclusively establishes, that the assertion or protected speech or petition activity was illegal as a matter of law, the defendant is precluded from using the anti-SLAPP statute to strike the plaintiff's action.'' Flatley's categorical-illegality principle extends to the scope of the Civil Code \textsection{} 47(b) litigation privilege: conduct that is criminal or illegal as a matter of law is not shielded by \textsection{} 47(b). To the extent the February 25, 2021 denial letter constituted a knowing misrepresentation of a material fact --- namely, that CSAA had ``concluded'' its ``investigation'' and had ``carefully reviewed the facts and circumstances'' when foundational first-responder records had not been obtained --- that representation constitutes fraud within the meaning of Civil Code \textsection{}\textsection{} 1709, 1710 and, under the analytical framework of Flatley, is categorically outside the scope of \textsection{} 47(b). This carve-out is independently sufficient to defeat any litigation-privilege defense regardless of whether the Action Apartment pre-litigation factors in ¶ 16A are also resolved in Plaintiff's favor.

\begin{center}\rule{0.5\linewidth}{0.5pt}\end{center}

\phantomsection\label{toc:section-iii-attorneys-fees-cure}
\section{SECTION III: ATTORNEY'S FEES CURE}\label{section-iii-attorneys-fees-cure}

\subsection{Deletion of Pro-Per Fee Demands}\label{deletion-of-pro-per-fee-demands}

\begin{enumerate}
\def\labelenumi{\arabic{enumi}.}
\setcounter{enumi}{21}
\tightlist
\item
  Plaintiff, proceeding in pro per, does not seek attorney's fees for self-representation in this action. \emph{Trope v. Katz} (1995) 11 Cal.4th 274.
\end{enumerate}

\subsection{Prentice Fees as Damages Element}\label{prentice-fees-as-damages-element}

\begin{enumerate}
\def\labelenumi{\arabic{enumi}.}
\setcounter{enumi}{22}
\item
  Under \emph{Prentice v. North American Title Guaranty Corp.} (1963) 59 Cal.2d 618, 620, a plaintiff who has been induced by fraud to enter into litigation may recover, as an element of tort damages, the attorney's fees and litigation expenses incurred in that litigation.
\item
  Plaintiff alleges that:
\end{enumerate}

\textbf{(Element 1 --- Fraud):} Defendant committed fraud in denying Plaintiff's claim and misrepresenting the basis for the denial;

\textbf{(Element 2 --- Induced Litigation):} As a proximate result of Defendant's fraud, Plaintiff was induced to incur attorney's fees in litigating CGC-21-594102; and

\textbf{(Element 3 --- Recoverable Fees):} The attorney's fees incurred in CGC-21-594102 are recoverable as an element of tort damages, not as a fee award.

24A. \textbf{Prentice Line-Item Enumeration and Induced-Litigation Element.} For purposes of the \emph{Prentice} three-element test pleaded in ¶ 24, Plaintiff alleges that but for CSAA's February 25, 2021 misrepresentation of investigative completeness --- which falsely conveyed to Plaintiff, an unrepresented direct third-party claimant, that liability had been investigated and rejected on a ``concluded investigation'' --- Plaintiff either would not have filed CGC-21-594102 at all, or would have filed it on materially different terms (e.g., with counsel retained under different terms, with a different scope of discovery, or within a different scheduling posture). That induced-litigation element is pleaded as a specific causal link between the fraud and the separate tort-measure damages recoverable under \emph{Prentice v. North American Title Guaranty Corp.} (1963) 59 Cal.2d 618, 620--21. Plaintiff seeks \emph{Prentice} damages in at least the following four line-item categories, each pleaded as tort damages and not as a prevailing-party fee award:

\textbf{(Line-Item 1 --- Professional Fees and Costs Incurred in CGC-21-594102):} All attorney's fees, co-counsel fees, consulting-counsel fees, and paraprofessional fees actually incurred in the underlying action; and all expert-witness fees, deposition fees, court-reporter fees, transcript costs, filing fees, motion-fees, subpoena and process-service fees, mediation fees, and related out-of-pocket litigation expenses actually incurred in that action.

\textbf{(Line-Item 2 --- The \$75,000 Cost Judgment as Induced-Litigation Damages):} To the extent the trial court in CGC-21-594102 has entered, or may enter, a cost award or cost judgment against Plaintiff in favor of the underlying defendant (CSAA's insured) in an approximate amount of \$75,000, that cost judgment is itself a \emph{Prentice} damages element: it is a sum of money Plaintiff has been compelled to bear by reason of having been fraudulently induced to file and prosecute CGC-21-594102 in response to CSAA's February 25, 2021 misrepresentation. Plaintiff pleads the \$75,000 cost judgment (or such final amount as is entered) as tort damages recoverable against CSAA under \emph{Prentice}, not as a re-litigation of the cost-judgment itself in CGC-21-594102. This pleading converts the defense's own MPA2 footnote 2 rhetoric --- that Plaintiff ``owes approximately \$75,000.00 in costs'' --- into an affirmative \emph{Prentice} damages item, because those very costs were incurred as a direct consequence of the fraud now in suit.

\textbf{(Line-Item 3 --- Transcript, Record, and Appellate Preparation Costs):} Costs of preparing, ordering, and certifying trial transcripts (including the October 29, 2024 Day-7 transcript), pretrial transcripts, clerk's transcripts, register-of-actions copies, and other record-on-appeal materials to the extent incurred in connection with the underlying action.

\textbf{(Line-Item 4 --- Investigative and Records-Acquisition Costs):} Costs of obtaining the very 911 audio, CAD records, BWC footage, and related first-responder records that CSAA itself failed to obtain before the February 25, 2021 denial --- each of which Plaintiff has had to procure (or is in the process of procuring) through public-records requests, subpoenas, or other third-party channels in order to litigate the fraud that would not have existed had CSAA performed the investigation it represented.

Line-Items 1--4 are pleaded as damages elements, not as a fee shift, and are not duplicative of any cost award under CCP \textsection{} 1032 in this action.

\begin{center}\rule{0.5\linewidth}{0.5pt}\end{center}

\phantomsection\label{toc:section-iv-punitive-damages-specificity}
\section{SECTION IV: PUNITIVE DAMAGES SPECIFICITY}\label{section-iv-punitive-damages-specificity}

\subsection{Particulars Block: Who, What, When, Where, How, By What Means}\label{particulars-block-who-what-when-where-how-by-what-means}

\begin{enumerate}
\def\labelenumi{\arabic{enumi}.}
\setcounter{enumi}{24}
\tightlist
\item
  Plaintiff seeks punitive damages under Civil Code \textsection{} 3294. In compliance with the particularity requirements of \emph{Lazar v. Superior Court} (1996) 12 Cal.4th 631, Plaintiff alleges the following specifics:
\end{enumerate}

\textbf{(a) WHO:} The February 25, 2021 denial letter was signed by \textbf{Sarah Rash, Claims Representative}, employed by CSAA Insurance Exchange. Sarah Rash acted within the scope of employment and under the direction and supervision of CSAA management.

\textbf{(b) WHAT:} The denial letter falsely represented that CSAA had ``concluded our investigation'' and had ``carefully review{[}ed{]} the facts'' when CSAA had not conducted the investigation required by Title 10 CCR \textsection{} 2695.7(b). The foundational incident records---911 audio, BWC footage, CAD logs---had not been obtained or reviewed.

\textbf{(c) WHEN:} The denial letter was dated \textbf{February 25, 2021}---just 21 days after the collision. The fraudulent claims handling process began on or about February 4, 2021 and continued through the date of the denial.

\textbf{(d) WHERE:} The claims handling occurred at CSAA's claims office. Communications were directed to Plaintiff at Plaintiff's residence in San Francisco, California.

\textbf{(e) HOW:} CSAA's claims representative prepared and sent the denial letter without conducting the investigation required by Title 10 CCR. The representative knew or should have known that the denial was not supported by a completed investigation.

\textbf{(f) BY WHAT MEANS:} The denial was communicated by written letter sent via U.S. Mail to Plaintiff.

25A. \textbf{Pleading-Stage Standard vs.~Trial-Stage Evidentiary Standard --- G.D. Searle Bright Line.} The defense is anticipated to cite \emph{In re Angelia P.} (1981) 28 Cal.3d 908 for the proposition that ``clear and convincing'' evidence is required at the pleading stage. Plaintiff affirmatively pleads that Angelia P. describes a \textbf{trial-stage} burden of proof on punitive damages and does not purport to govern the \textbf{pleading} sufficiency of a \textsection{} 3294 allegation. Under \emph{G.D. Searle \& Co.~v. Superior Court} (1975) 49 Cal.App.3d 22, 29, ``the inquiry on {[}demurrer or strike{]} is simply whether the complaint alleges facts which, if proven by clear and convincing evidence at trial, would support a claim for punitive damages.'' The pleading stage therefore requires only that the complaint plead facts --- with \emph{Lazar v. Superior Court} (1996) 12 Cal.4th 631 particularity --- which, if proven by clear and convincing evidence at trial, would authorize a jury to find malice, oppression, or fraud under Civil Code \textsection{} 3294(c). The particularized block at ¶ 25(a)--(f), the managing-agent / ratification block at ¶¶ 26--27, the deny-delay-defend pattern at ¶ 28, and the confession/victim/eyewitness triple at ¶¶ 29--30 together satisfy the \emph{G.D. Searle} pleading test. Whether the ultimate trier of fact will find the allegations proven by clear and convincing evidence is a trial question, not a \textsection{} 436 strike or \textsection{} 430.10 demurrer question.

\subsection{Managing-Agent / Ratification Paragraph Under \textsection{} 3294(b)}\label{managing-agent-ratification-paragraph-under-3294b}

\begin{enumerate}
\def\labelenumi{\arabic{enumi}.}
\setcounter{enumi}{25}
\item
  Civil Code \textsection{} 3294(b) requires that, to recover punitive damages against a corporate defendant, plaintiff must prove that an officer, director, or managing agent of the corporation committed, authorized, or ratified the wrongful conduct.
\item
  Plaintiff alleges on information and belief that:
\end{enumerate}

\begin{enumerate}
\def\labelenumi{(\alph{enumi})}
\item
  Sarah Rash was acting under the direct supervision of a managing agent of CSAA Insurance Exchange;
\item
  The denial was reviewed and approved by a supervisor or manager with authority over claims decisions;
\item
  CSAA's claim-denial-approval policy, which resulted in the inadequate investigation and wrongful denial, was established or ratified by officers, directors, or managing agents; and
\item
  After learning of the wrongful denial (through this litigation or otherwise), CSAA's managing agents ratified the conduct by failing to correct the denial.
\end{enumerate}

\subsection{``Deny-Delay-Defend'' Institutional Pattern}\label{deny-delay-defend-institutional-pattern}

\begin{enumerate}
\def\labelenumi{\arabic{enumi}.}
\setcounter{enumi}{27}
\item
  Plaintiff alleges that CSAA's failure to investigate was not inadvertent or a product of ordinary claims-handling error. CSAA institutionalized a ``deny-delay-defend'' practice that systematically omits foundational investigative steps, then conceals the omission through denial letters representing the investigation as concluded.
\item
  CSAA's concealment is particularly egregious because the concealed evidence was not ambiguous or contested, and each source was identifiable, available, and obtainable by CSAA through ordinary claim-handling channels within days of the February 4, 2021 collision:
\end{enumerate}

\begin{enumerate}
\def\labelenumi{(\alph{enumi})}
\item
  The \textbf{911 audio recording} --- associated with SFPD Incident Report No.~210083458 (conformed to proof) and the corresponding San Francisco Department of Emergency Management CAD event on February 4, 2021 --- contained CSAA's insured's own acknowledgment of contact with Plaintiff;
\item
  The \textbf{Body-Worn Camera (BWC) footage} from the responding SFPD officers (names and star numbers in SFPD's custody, available through a Penal Code \textsection{} 832.7(b) public-records request and/or subpoena) contained an independent eyewitness, \textbf{Justin Rosemond}, stating on-camera that the driver struck Plaintiff; and
\item
  Plaintiff's own \textbf{contemporaneous 911 report}, made on February 4, 2021 from the collision scene and reflected in both the SFPD CAD record and the SFPD Incident Report, stated that the driver hit him.
\end{enumerate}

Each of the three sources (a)--(c) was identifiable to CSAA from the face of the SFPD Incident Report, each was available to CSAA by routine public-records request or subpoena within days of the collision, and each was available before the February 25, 2021 denial.

\begin{enumerate}
\def\labelenumi{\arabic{enumi}.}
\setcounter{enumi}{29}
\tightlist
\item
  CSAA denied a claim supported by \textbf{confession, victim report, and eyewitness corroboration}. This conduct constitutes ``malice'' and ``despicable conduct'' within the meaning of Civil Code \textsection{} 3294 and the \emph{American Airlines, Inc.~v. Sheppard, Mullin, Richter \& Hampton} (2018) 28 Cal.App.5th 1077 standard.
\end{enumerate}

\begin{center}\rule{0.5\linewidth}{0.5pt}\end{center}

\phantomsection\label{toc:section-v-damages}
\section{SECTION V: DAMAGES}\label{section-v-damages}

\subsection{(a) Reliance Damages --- Civil Code \textsection{} 1709}\label{a-reliance-damages-civil-code-1709}

\begin{enumerate}
\def\labelenumi{\arabic{enumi}.}
\setcounter{enumi}{30}
\tightlist
\item
  Plaintiff seeks reliance damages under Civil Code \textsection{} 1709 for Defendant's fraud. Specifically, Plaintiff incurred the following expenses in reasonable reliance on Defendant's misrepresentation of investigative completeness:
\end{enumerate}

\begin{enumerate}
\def\labelenumi{(\roman{enumi})}
\tightlist
\item
  Costs of investigating and challenging the denial;
\item
  Filing fees and court costs in CGC-21-594102;
\item
  Expert witness fees; and
\item
  Other out-of-pocket expenses incurred in reasonable reliance on Defendant's representations.
\end{enumerate}

\subsection{(b) Prentice Fees}\label{b-prentice-fees}

\begin{enumerate}
\def\labelenumi{\arabic{enumi}.}
\setcounter{enumi}{31}
\tightlist
\item
  Plaintiff seeks attorney's fees as tort damages under \emph{Prentice v. North American Title Guaranty Corp.} (1963) 59 Cal.2d 618, as set forth in Section III above.
\end{enumerate}

\subsection{(c) Lost Earnings --- Fraud-Delta Only}\label{c-lost-earnings-fraud-delta-only}

\begin{enumerate}
\def\labelenumi{\arabic{enumi}.}
\setcounter{enumi}{32}
\tightlist
\item
  Plaintiff seeks lost earnings damages, but only to the extent such losses were directly caused by Defendant's fraud (``fraud-delta''). This claim is limited to the differential between what Plaintiff earned and what Plaintiff would have earned but for Defendant's fraudulent concealment of the non-investigation.
\end{enumerate}

\subsection{(d) Physical Exacerbation}\label{d-physical-exacerbation}

\begin{enumerate}
\def\labelenumi{\arabic{enumi}.}
\setcounter{enumi}{33}
\tightlist
\item
  Plaintiff alleges that Defendant's fraud exacerbated Plaintiff's pre-existing physical condition:
\end{enumerate}

\begin{enumerate}
\def\labelenumi{(\alph{enumi})}
\tightlist
\item
  \textbf{Foreseeability:} It was foreseeable that wrongful denial of a claim would cause increased stress and physical symptoms;
\item
  \textbf{Eggshell Plaintiff:} Defendant takes Plaintiff as Defendant finds him;
\item
  \textbf{Objectively Billed:} Plaintiff seeks only medical expenses actually billed and causally linked to the exacerbation.
\end{enumerate}

\subsection{(e) Emotional Distress --- Parasitic to Fraud Only}\label{e-emotional-distress-parasitic-to-fraud-only}

\begin{enumerate}
\def\labelenumi{\arabic{enumi}.}
\setcounter{enumi}{34}
\tightlist
\item
  Plaintiff seeks emotional distress damages, but only as parasitic damages accompanying the fraud tort. Plaintiff does not assert a standalone claim for intentional infliction of emotional distress.
\end{enumerate}

\subsection{(f) Special Damages}\label{f-special-damages}

\begin{enumerate}
\def\labelenumi{\arabic{enumi}.}
\setcounter{enumi}{35}
\tightlist
\item
  Plaintiff seeks special damages for fraud-caused out-of-pocket expenses, including:
\end{enumerate}

\begin{enumerate}
\def\labelenumi{(\alph{enumi})}
\tightlist
\item
  Transcript, court-reporter, and certified-copy costs;
\item
  Subpoena, process-service, and certification fees for obtaining 911, BWC, and CAD materials;
\item
  Document-preparation, authentication, and verification costs;
\item
  Professional and logistical costs of reconstructing the claim history.
\end{enumerate}

\subsection{(g) Future Damages --- CACI 3903A Foundation}\label{g-future-damages-caci-3903a-foundation}

\begin{enumerate}
\def\labelenumi{\arabic{enumi}.}
\setcounter{enumi}{36}
\tightlist
\item
  To the extent Plaintiff seeks future damages, Plaintiff alleges a foundation under CACI 3903A: Plaintiff's injuries are ongoing; it is reasonably certain that Plaintiff will incur future medical expenses and/or lost earnings; and the amount of future damages can be estimated with reasonable certainty.
\end{enumerate}

\subsection{(h) UCL Injunction --- Specific Misrepresentation Pattern}\label{h-ucl-injunction-specific-misrepresentation-pattern}

\begin{enumerate}
\def\labelenumi{\arabic{enumi}.}
\setcounter{enumi}{37}
\tightlist
\item
  Plaintiff seeks injunctive relief under Business and Professions Code \textsection{} 17203 to:
\end{enumerate}

\begin{enumerate}
\def\labelenumi{(\alph{enumi})}
\item
  Enjoin Defendant from \textbf{issuing denial letters stating an investigation is ``concluded''} when foundational incident records have not been obtained or reviewed;
\item
  Require Defendant to comply with Title 10 CCR \textsection{}\textsection{} 2695.4 and 2695.7 by conducting thorough investigations before denying claims; and
\item
  Require Defendant to maintain adequate documentation of claim investigations.
\end{enumerate}

\subsection{(i) Costs}\label{i-costs}

\begin{enumerate}
\def\labelenumi{\arabic{enumi}.}
\setcounter{enumi}{38}
\tightlist
\item
  Plaintiff seeks costs of suit under Code of Civil Procedure \textsection{}\textsection{} 1032 and 1033.5 for costs incurred in this action (CGC-25-631802 only).
\end{enumerate}

\subsection{(j) Prejudgment Interest}\label{j-prejudgment-interest}

\begin{enumerate}
\def\labelenumi{\arabic{enumi}.}
\setcounter{enumi}{39}
\tightlist
\item
  Plaintiff seeks prejudgment interest under Civil Code \textsection{}\textsection{} 3287(a) and 3288.
\end{enumerate}

\begin{center}\rule{0.5\linewidth}{0.5pt}\end{center}

\phantomsection\label{toc:section-vi-primary-rights-reservation}
\section{SECTION VI: PRIMARY RIGHTS / RESERVATION}\label{section-vi-primary-rights-reservation}

\subsection{Distinct Primary Right}\label{distinct-primary-right}

\begin{enumerate}
\def\labelenumi{\arabic{enumi}.}
\setcounter{enumi}{40}
\tightlist
\item
  The primary right asserted here is the right to be free from deliberate misrepresentation by a commercial actor regarding the completeness of an investigative process that directly determined Plaintiff's economic options as a claimant. The invasion of that right occurred in February 2021 when CSAA caused Plaintiff to believe an investigation had concluded when it had not, thereby destroying Plaintiff's ability to independently evaluate the true state of the claim and make informed decisions about preservation, evidence-gathering, counsel retention, and negotiation strategy.
\end{enumerate}

41A. \textbf{Lucido Five-Factor Walk-Through --- Primary-Rights and Collateral-Estoppel Preemption.} To preempt any anticipated defense argument that CGC-21-594102 (the underlying personal-injury negligence action against CSAA's insured) collaterally estops or otherwise precludes this action, Plaintiff pleads the five-factor issue-preclusion test of \emph{Lucido v. Superior Court} (1990) 51 Cal.3d 335, 341 against the CGC-21-594102 record: \textbf{(i) Identical Issue:} The issue decided in CGC-21-594102 was whether Subhi Abdelhalim breached a Vehicle Code duty of care toward Plaintiff in the February 4, 2021 collision. The issue in this action is whether CSAA Insurance Exchange committed common law fraud and engaged in unfair business practices in its February 25, 2021 claim-handling representation that its investigation had ``concluded.'' These are not identical issues. \textbf{(ii) Actually Litigated:} Whether CSAA's 2021 claim-handling constituted fraud was not litigated in CGC-21-594102; CSAA was not a party to that action and no fraud claim was tendered to the jury. \textbf{(iii) Necessarily Decided:} The CGC-21-594102 9--3 negligence verdict did not necessarily decide --- and could not have necessarily decided --- any issue of CSAA's claim-handling fraud, because that issue was not before the trier of fact. \textbf{(iv) Final Judgment:} The CGC-21-594102 judgment is presently on appeal (A173827), and a judgment on appeal is not final for preclusion purposes. \textbf{(v) Same Party or Privy:} CSAA was not a named party in CGC-21-594102 and its privity with the insured, if any, runs only to the insured's negligence liability, not to CSAA's independent claim-handling conduct toward a direct third-party claimant. Under the primary-rights doctrine of \emph{Mycogen Corp.~v. Monsanto Co.} (2002) 28 Cal.4th 888, 904, and \emph{Crowley v. Katleman} (1994) 8 Cal.4th 666, 681--82, this action and CGC-21-594102 arise from different primary rights --- the right to be free from a commercial fraud versus the right to be free from negligent driving --- and therefore do not share a single ``cause of action'' for preclusion purposes. Each of the five \emph{Lucido} factors independently defeats issue preclusion, and the primary-rights framework defeats claim preclusion.

\begin{enumerate}
\def\labelenumi{\arabic{enumi}.}
\setcounter{enumi}{41}
\tightlist
\item
  This primary right is distinct from:
\end{enumerate}

\begin{enumerate}
\def\labelenumi{(\alph{enumi})}
\tightlist
\item
  The personal injury primary right litigated in CGC-21-594102;
\item
  The extrinsic fraud primary right asserted in CGC-25-631801; and
\item
  Any right to policy benefits (not asserted here).
\end{enumerate}

\subsection{Express Disclaimer of PI Damages}\label{express-disclaimer-of-pi-damages}

\begin{enumerate}
\def\labelenumi{\arabic{enumi}.}
\setcounter{enumi}{42}
\tightlist
\item
  Plaintiff expressly disclaims and does not seek in this action any personal injury damages that were or could have been awarded in CGC-21-594102.
\end{enumerate}

\subsection{2021-Forward Damages Anchor}\label{forward-damages-anchor}

\begin{enumerate}
\def\labelenumi{\arabic{enumi}.}
\setcounter{enumi}{43}
\tightlist
\item
  All damages sought in this action arise from Defendant's conduct on or after February 25, 2021 (the date of the denial letter). Plaintiff does not seek damages for any conduct predating the February 25, 2021 denial.
\end{enumerate}

44A. \textbf{Express Double-Recovery Reservation Between CGC-25-631802 and CGC-25-631801.} Plaintiff is simultaneously prosecuting CGC-25-631801, an independent action in equity for extrinsic fraud (vacatur of the CGC-21-594102 judgment procured by fraud on the court). The theory of this action (CGC-25-631802) is distinct from, and non-duplicative of, the equitable-vacatur theory in CGC-25-631801: this action seeks legal-damages relief against CSAA for its 2021 claim-handling fraud, while CGC-25-631801 seeks equitable relief against CSAA and defense counsel for concealment of evidence during the 2024--2025 trial of CGC-21-594102. To prevent any double recovery and to foreclose any election-of-remedies argument, Plaintiff expressly: (a) disclaims, in this action, any damages measure that would also be recovered as equitable or restitutionary relief under CGC-25-631801 (including any restitution-of-the-\$75,000-cost-judgment that would duplicate equitable vacatur relief in CGC-25-631801); (b) reserves the right, pre-judgment, to elect remedies between the two actions upon any ruling that places them in tension; and (c) affirms that the \emph{Prentice} line-items in ¶ 24A are sought here only to the extent they are not awarded through equitable vacatur in CGC-25-631801. Any overlap will be managed by pre-judgment election under the principles of \emph{Mycogen Corp.~v. Monsanto Co.} (2002) 28 Cal.4th 888, 904--05, and \emph{Crowley v. Katleman} (1994) 8 Cal.4th 666, to ensure that Plaintiff recovers once and only once for any single item of loss.

\begin{center}\rule{0.5\linewidth}{0.5pt}\end{center}

\phantomsection\label{toc:section-vii-title-10-investigation-standard}
\section{SECTION VII: TITLE 10 INVESTIGATION STANDARD}\label{section-vii-title-10-investigation-standard}

\subsection{What Thorough Investigation Requires}\label{what-thorough-investigation-requires}

\begin{enumerate}
\def\labelenumi{\arabic{enumi}.}
\setcounter{enumi}{44}
\tightlist
\item
  Under Title 10 CCR \textsection{} 2695.7(b), a thorough investigation of an insurance claim requires, at minimum:
\end{enumerate}

\begin{enumerate}
\def\labelenumi{(\alph{enumi})}
\tightlist
\item
  Reviewing all information submitted by the claimant;
\item
  Obtaining and reviewing relevant third-party evidence, including police reports, medical records, and witness statements;
\item
  Obtaining available public records, including 911 recordings, CAD records, and body-worn camera footage where relevant;
\item
  Documenting the investigation in the claim file; and
\item
  Basing any denial on a reasonable evaluation of all available evidence.
\end{enumerate}

\subsection{What CSAA Did Not Do}\label{what-csaa-did-not-do}

\begin{enumerate}
\def\labelenumi{\arabic{enumi}.}
\setcounter{enumi}{45}
\tightlist
\item
  Defendant failed to conduct a thorough investigation as required by Title 10 CCR. Specifically:
\end{enumerate}

\begin{enumerate}
\def\labelenumi{(\alph{enumi})}
\tightlist
\item
  Defendant did not obtain or review the 911 recordings from the date of the collision;
\item
  Defendant did not obtain or review body-worn camera footage from responding officers;
\item
  Defendant did not obtain or review CAD records documenting the emergency response;
\item
  Defendant did not interview witnesses whose contact information was available;
\item
  Defendant did not adequately document the basis for its denial in the claim file; and
\item
  Defendant issued a denial based on an incomplete and one-sided review of the evidence.
\end{enumerate}

\subsection{Clarification: Not a Private Ins. Code \textsection{} 790.03 Claim}\label{clarification-not-a-private-ins.-code-790.03-claim}

\begin{enumerate}
\def\labelenumi{\arabic{enumi}.}
\setcounter{enumi}{46}
\tightlist
\item
  This section supports the UCL ``unlawful'' prong and is \textbf{not} pleaded as a private Insurance Code \textsection{} 790.03 cause of action. \emph{Moradi-Shalal v. Fireman's Fund Insurance Companies} (1988) 46 Cal.3d 287. The \textsection{} 790.03 violations are predicate conduct for UCL purposes only. \emph{State Farm Fire \& Casualty Co.~v. Superior Court} (1996) 45 Cal.App.4th 1093, 1103--04.
\end{enumerate}

47A. \textbf{Insurance Code \textsection{} 790.03(h)(5) --- Liability-Reasonably-Clear Standard as UCL Unlawful Predicate.} In addition to the Title 10 CCR \textsection{} 2695.4 and \textsection{} 2695.7 predicates pleaded at ¶¶ 45--46, Plaintiff alleges that CSAA's February 25, 2021 denial violated Insurance Code \textsection{} 790.03(h)(5), which defines as an unfair claims settlement practice ``{[}n{]}ot attempting in good faith to effectuate prompt, fair, and equitable settlements of claims in which liability has become reasonably clear.'' Liability was reasonably clear to CSAA on or before February 25, 2021 because each of the three foundational first-responder sources (911 audio / CAD, SFPD BWC footage with on-camera independent-eyewitness corroboration from Justin Rosemond, and Plaintiff's contemporaneous 911 report) was identifiable from the face of the SFPD Incident Report (Report No.~210083458, conformed to proof) and available to CSAA by routine public-records request or subpoena. Pleading \textsection{} 790.03(h)(5) here is not an assertion of a private right of action under \textsection{} 790.03 --- which remains precluded by \emph{Moradi-Shalal v. Fireman's Fund Insurance Cos.} (1988) 46 Cal.3d 287 --- but is pleaded \textbf{solely as a UCL unlawful predicate} consistent with the \emph{Moradi-Shalal} carve-out recognized in \emph{State Farm Fire \& Casualty Co.~v. Superior Court} (1996) 45 Cal.App.4th 1093, 1103--04. This paragraph is an independent UCL unlawful-prong predicate, in addition to the common-law-fraud-based unlawful-prong predicate in the Third Cause of Action at ¶ 65.

\begin{center}\rule{0.5\linewidth}{0.5pt}\end{center}

\phantomsection\label{toc:section-viii-delayed-discovery}
\section{SECTION VIII: DELAYED DISCOVERY}\label{section-viii-delayed-discovery}

\subsection{Three-Anchor Discovery Accrual}\label{three-anchor-discovery-accrual}

\begin{enumerate}
\def\labelenumi{\arabic{enumi}.}
\setcounter{enumi}{47}
\item
  Plaintiff did not discover the facts constituting the fraud at the time of the February 25, 2021 denial, as the denial letter actively concealed the non-investigation.
\item
  The discovery accrual is anchored to three consecutive record events:
\end{enumerate}

\begin{enumerate}
\def\labelenumi{(\alph{enumi})}
\item
  \textbf{January 17, 2025:} Defense MIL No.~17 filings asserted the 911 and related materials had not been produced in discovery and were of ``unknown provenance''---revealing that CSAA's 2021 representation of investigative completeness was false;
\item
  \textbf{February 18, 2025:} Pretrial hearing testimony by Jeremy Jessup stated those materials had been produced to defense counsel in August-September 2022, while defense counsel simultaneously stated the 911 call was ``news to me,'' creating a direct internal contradiction; and
\item
  \textbf{April 16, 2025:} Trial admission by defense counsel that the materials ``were provided by Dolan,'' confirming the materials had been in defense possession all along.
\end{enumerate}

\begin{enumerate}
\def\labelenumi{\arabic{enumi}.}
\setcounter{enumi}{49}
\item
  These three events, taken together, constitute the first point at which Plaintiff could reasonably have determined that CSAA's 2021 representation of investigative completeness was false.
\item
  Plaintiff filed this action within the applicable statute of limitations following these discovery events.
\end{enumerate}

\begin{center}\rule{0.5\linewidth}{0.5pt}\end{center}

\phantomsection\label{toc:causes-of-action}
\section{CAUSES OF ACTION}\label{causes-of-action}

\subsection{FIRST CAUSE OF ACTION: FRAUD --- INTENTIONAL OMISSION AND EXTRINSIC FRAUD}\label{first-cause-of-action-fraud-intentional-omission-and-extrinsic-fraud}

(Against CSAA Insurance Exchange and DOES 1 through 50, inclusive)

\begin{enumerate}
\def\labelenumi{\arabic{enumi}.}
\setcounter{enumi}{51}
\tightlist
\item
  Plaintiff incorporates by reference all preceding paragraphs.
\end{enumerate}

52A. \textbf{Wrongful-Conduct Scope --- Express Anti-SLAPP Disavowal.} The wrongful conduct pleaded in this cause of action is \textbf{not} predicated on, and Plaintiff expressly disavows any theory of liability predicated on: (a) the filing or prosecution of CGC-21-594102 by CSAA's insured or by defense counsel; (b) any motion, in-limine application, pleading, or other court filing made in connection with CGC-21-594102; (c) statements or arguments made in open court, in depositions, or in written discovery responses in CGC-21-594102; or (d) any other act or statement that would constitute ``an act in furtherance of a person's right of petition or free speech . . . in connection with a public issue'' within the meaning of Code of Civil Procedure \textsection{} 425.16(e). The wrongful conduct in this cause of action is CSAA's pre-litigation, non-communicative 2021 claim-handling conduct --- the failure to investigate, the issuance of a denial letter misrepresenting investigative completeness, and the omission of foundational first-responder evidence --- which occurred before any litigation was pending or contemplated and which, under ¶¶ 16A and 21A above, is categorically outside the protected-activity scope of Civil Code \textsection{} 47(b) and Code of Civil Procedure \textsection{} 425.16. This paragraph eliminates any overlap between CSAA's anticipated anti-SLAPP motion and the CCP \textsection{} 436 motion to strike, because no element of this cause of action relies on litigation-conduct activity protected by \textsection{} 425.16 step 1.

\begin{enumerate}
\def\labelenumi{\arabic{enumi}.}
\setcounter{enumi}{52}
\item
  The wrongful conduct pleaded in this cause of action is CSAA's non-communicative failure to investigate the claim, its intentional omission of obtaining material evidence, and its active concealment of that failure (extrinsic fraud)---not any later petitioning conduct, motion practice, or in-court advocacy.
\item
  CSAA had a duty to investigate the claim fairly, yet it failed to do so. It actively concealed this failure by providing Plaintiff with a denial letter containing the false representation that it had ``carefully reviewed the facts.''
\item
  The denial letter serves as evidence of the omission and the extrinsic fraud, showing that CSAA intended Plaintiff to rely on the false completeness to suppress his claim.
\item
  Plaintiff justifiably relied on the concealment.
\item
  Plaintiff suffered damage as a proximate result.
\end{enumerate}

\subsection{SECOND CAUSE OF ACTION: FRAUD --- CONCEALMENT / HALF-TRUTH}\label{second-cause-of-action-fraud-concealment-half-truth}

(Against CSAA Insurance Exchange and DOES 1 through 50, inclusive)

\begin{enumerate}
\def\labelenumi{\arabic{enumi}.}
\setcounter{enumi}{57}
\item
  Plaintiff incorporates by reference all preceding paragraphs.
\item
  The wrongful conduct pleaded in this cause of action is the non-communicative institutional failure to investigate, which CSAA concealed via misleading half-truths in its ordinary business communications.
\item
  By issuing a denial ``based on the facts currently available,'' CSAA omitted the material fact that it had not yet obtained, reviewed, or responsibly accounted for foundational incident information---specifically, the 911 audio, BWC footage, and CAD logs.
\item
  Plaintiff reasonably relied on the resulting half-truth and omission.
\item
  Plaintiff suffered damages as a proximate result.
\end{enumerate}

\subsection{THIRD CAUSE OF ACTION: VIOLATION OF BUSINESS AND PROFESSIONS CODE \textsection{} 17200 ET SEQ.}\label{third-cause-of-action-violation-of-business-and-professions-code-17200-et-seq.}

(Against CSAA Insurance Exchange and DOES 1 through 50, inclusive)

\begin{enumerate}
\def\labelenumi{\arabic{enumi}.}
\setcounter{enumi}{62}
\item
  Plaintiff incorporates by reference all preceding paragraphs.
\item
  Plaintiff pleads this cause of action in the alternative and as subordinate to the fraud causes of action.
\item
  \textbf{(a) Unlawful prong:} CSAA's conduct is unlawful because it constitutes common law fraud under Civil Code \textsection{}\textsection{} 1709--1710, independently pleaded in the First and Second Causes of Action above. This UCL theory does not depend on, and expressly disclaims, any private enforcement of Insurance Code \textsection{} 790.03(h).
\item
  \textbf{(b) Fraudulent prong:} CSAA's systematic misrepresentation of investigation completeness to unrepresented third-party claimants constitutes a business practice likely to deceive the public. The representation that a ``careful investigation'' has ``concluded'' when foundational incident records have not been obtained or reviewed is a pattern of institutional conduct that deceives a class of claimants as a matter of practice.
\end{enumerate}

66A. \textbf{Reasonable-Consumer Test --- UCL Fraudulent-Prong Standard.} The UCL fraudulent prong is governed by the ``reasonable consumer'' test: a business practice is ``fraudulent'' within the meaning of Business and Professions Code \textsection{} 17200 if ``members of the public are likely to be deceived'' by the practice. \emph{Williams v. Gerber Products Co.} (9th Cir. 2008) 552 F.3d 934, 938; \emph{Chapman v. Skype, Inc.} (2013) 220 Cal.App.4th 217, 226. Plaintiff alleges that a reasonable unrepresented third-party claimant in CSAA's claims-intake universe is likely to be deceived by a denial letter that represents the insurer's investigation as ``concluded'' after ``careful review'' when foundational first-responder records (911 audio, CAD, BWC, incident reports) have not been obtained, because a reasonable claimant, lacking independent means to audit an insurer's internal claims file, will rely on the insurer's representation of investigative completeness as a basis for assessing his or her claim and its viability. The reasonable-consumer test is a pleading-stage question where the pleaded facts plausibly support likely deception; whether the ultimate trier of fact will find the practice deceptive by a preponderance of evidence is a trial question.

\begin{enumerate}
\def\labelenumi{\arabic{enumi}.}
\setcounter{enumi}{66}
\item
  \textbf{(c) Unfair prong:} The ``deny-delay-defend'' scheme constitutes an unfair business practice because it produces a commercial benefit to CSAA (suppression of valid claims, reduction of investigation costs, reserve protection) through means that cause harm to claimants disproportionate to any legitimate claims-management justification.
\item
  Plaintiff alleges he lost money or property as a result of that conduct, including documentable investigation, verification, record-obtaining, and related out-of-pocket expenditures.
\item
  Plaintiff seeks under this cause of action injunctive relief and, only to the extent legally available, restitutionary relief.
\end{enumerate}

\begin{center}\rule{0.5\linewidth}{0.5pt}\end{center}

\phantomsection\label{toc:prayer-for-relief}
\section{PRAYER FOR RELIEF}\label{prayer-for-relief}

WHEREFORE, Plaintiff prays for judgment against Defendants, and each of them, as follows:

\begin{enumerate}
\def\labelenumi{\arabic{enumi}.}
\item
  For general and special damages according to proof on the fraud causes of action;
\item
  For reliance damages under Civil Code \textsection{} 1709;
\item
  For \emph{Prentice} fees as tort damages;
\item
  For restitutionary and injunctive relief to the extent legally available under Business and Professions Code \textsection{} 17200 et seq., including an injunction prohibiting CSAA from issuing denial letters stating an investigation is ``concluded'' when foundational incident records have not been obtained or reviewed;
\item
  For punitive damages where legally recoverable and according to proof under Civil Code \textsection{} 3294;
\item
  For costs of suit incurred herein under CCP \textsection{}\textsection{} 1032, 1033.5;
\item
  For prejudgment interest as allowed by law under Civil Code \textsection{}\textsection{} 3287(a), 3288; and
\item
  For such other and further relief as the Court deems just and proper.
\end{enumerate}

\begin{center}\rule{0.5\linewidth}{0.5pt}\end{center}

\phantomsection\label{toc:demand-for-jury-trial}
\section{DEMAND FOR JURY TRIAL}\label{demand-for-jury-trial}

Plaintiff hereby demands trial by jury on all issues so triable.

\begin{center}\rule{0.5\linewidth}{0.5pt}\end{center}

\textbf{Dated:} \_\_\_\_\_\_\_\_\_\_\_\_\_\_\_

\textbf{FRANCISCUS DYLAN ROSARIO}\\
Plaintiff, In Pro Per\\
7624 Melody Drive\\
Rohnert Park, CA 94928\\
Tel: 650.488.7510\\
Email: soltrinox@gmail.com

\begin{center}\rule{0.5\linewidth}{0.5pt}\end{center}

\phantomsection\label{toc:verification}
\section{VERIFICATION}\label{verification}

I, FRANCISCUS DYLAN ROSARIO, am the Plaintiff in the above-entitled action. I have read the foregoing Third Amended Complaint and know its contents. The same is true of my own knowledge except as to matters stated on information and belief, and as to those matters I believe them to be true.

I declare under penalty of perjury under the laws of the State of California that the foregoing is true and correct.

Executed on \_\_\_\_\_\_\_\_\_\_\_\_\_\_\_ at \_\_\_\_\_\_\_\_\_\_\_\_\_\_\_, California.

\begin{center}\rule{0.5\linewidth}{0.5pt}\end{center}

FRANCISCUS DYLAN ROSARIO

\nolinenumbers
\vspace{6mm}
\noindent /s/ Franciscus Dylan Rosario\\
\vspace{4mm}
\noindent\includegraphics[height=0.5in]{../../../Support/signature.png}

\end{document}
